% ============================================================
%   LECTURE NOTES TEMPLATE
%   Supports: theorems, proofs, code, plots, geometry, figures
% ============================================================

\documentclass[11pt, a4paper]{article}

% ── Core packages ───────────────────────────────────────────
\usepackage[T1]{fontenc}
\usepackage[utf8]{inputenc}
\usepackage{lmodern}                   % Better fonts
\usepackage[margin=2.5cm]{geometry}    % Page margins
\usepackage{microtype}                 % Better text justification

% ── Math ────────────────────────────────────────────────────
\usepackage{amsmath, amssymb, amsthm, mathtools}
\usepackage{thmtools}                  % Extended theorem tools
\usepackage{bbm}                       % \mathbbm{1} indicator
\usepackage{bm}                        % Bold math \bm{}

% ── Colors ──────────────────────────────────────────────────
\usepackage[dvipsnames, svgnames, x11names]{xcolor}

% ── Define Color Palette ────────────────────────────────────
\definecolor{TheoremBG}    {RGB}{230, 245, 230}   % soft green
\definecolor{TheoremBorder}{RGB}{ 34, 139,  34}   % forest green
\definecolor{PropBG}       {RGB}{230, 240, 255}   % soft blue
\definecolor{PropBorder}   {RGB}{ 70, 130, 180}   % steel blue
\definecolor{CorolBG}      {RGB}{255, 245, 220}   % warm cream
\definecolor{CorolBorder}  {RGB}{205, 133,  63}   % peru
\definecolor{DefBG}        {RGB}{245, 230, 255}   % lavender
\definecolor{DefBorder}    {RGB}{128,   0, 128}   % purple
\definecolor{ExBG}         {RGB}{240, 248, 255}   % alice blue
\definecolor{ExBorder}     {RGB}{ 70, 130, 180}   % steel blue
\definecolor{RemBG}        {RGB}{255, 250, 240}   % floral white
\definecolor{RemBorder}    {RGB}{169, 169, 169}   % dark gray
\definecolor{LemmaBG}      {RGB}{230, 250, 250}   % light cyan
\definecolor{LemmaBorder}  {RGB}{ 32, 178, 170}   % light sea green
\definecolor{ProblemBG}    {RGB}{230, 240, 255}   % same as Prop
\definecolor{ProblemBorder}{RGB}{ 30,  90, 200}   % royal blue
\definecolor{CodeBG}       {RGB}{ 40,  44,  52}   % dark (One Dark)
\definecolor{CodeKey}      {RGB}{198, 120, 221}   % purple keyword
\definecolor{CodeStr}      {RGB}{152, 195, 121}   % green string
\definecolor{CodeCom}      {RGB}{ 92, 100, 112}   % gray comment
\definecolor{CodeNum}      {RGB}{209, 154, 102}   % orange number
\definecolor{LinkColor}    {RGB}{ 10,  80, 160}

% ── Theorem Environments (tcolorbox) ────────────────────────
\usepackage[most]{tcolorbox}
\tcbuselibrary{theorems, skins, breakable}

%  Macro for colored theorem boxes
\newcommand{\maketheorbox}[6]{%
  % #1 = env name, #2 = display name, #3 = BG color, #4 = border color
  % #5 = title text color, #6 = counter (share with: e.g. theorem)
  \tcbset{#1style/.style={%
    enhanced, breakable,
    colback=#3, colframe=#4,
    arc=4pt,
    fonttitle=\bfseries\color{white},
    attach boxed title to top left={yshift=-\tcboxedtitleheight/2, xshift=8pt},
    boxed title style={colback=#4, colframe=#4, sharp corners},
    sharp corners=south,
    boxrule=1pt,
    top=4mm, bottom=4mm,
  }}%
}

% ── Theorem / Proposition / etc. counters ───────────────────
\newcounter{thmcounter}[section]
\renewcommand{\thethmcounter}{\thesection.\arabic{thmcounter}}

% THEOREM
\newtcbtheorem[use counter=thmcounter]{theorem}{Theorem}{
  enhanced, breakable,
  colback=TheoremBG, colframe=TheoremBorder,
  arc=4pt,
  fonttitle=\bfseries\color{white},
  attach boxed title to top left={yshift=-\tcboxedtitleheight/2, xshift=8pt},
  boxed title style={colback=TheoremBorder, colframe=TheoremBorder,
    sharp corners, fontupper=\color{white}\bfseries},
  sharp corners=south, boxrule=1pt,
  top=5mm, bottom=4mm,
}{thm}

% PROPOSITION
\newtcbtheorem[use counter=thmcounter]{proposition}{Proposition}{
  enhanced, breakable,
  colback=PropBG, colframe=PropBorder,
  arc=4pt,
  fonttitle=\bfseries\color{white},
  attach boxed title to top left={yshift=-\tcboxedtitleheight/2, xshift=8pt},
  boxed title style={colback=PropBorder, colframe=PropBorder,
    sharp corners, fontupper=\color{white}\bfseries},
  sharp corners=south, boxrule=1pt,
  top=5mm, bottom=4mm,
}{prop}

% LEMMA
\newtcbtheorem[use counter=thmcounter]{lemma}{Lemma}{
  enhanced, breakable,
  colback=LemmaBG, colframe=LemmaBorder,
  arc=4pt,
  fonttitle=\bfseries\color{white},
  attach boxed title to top left={yshift=-\tcboxedtitleheight/2, xshift=8pt},
  boxed title style={colback=LemmaBorder, colframe=LemmaBorder,
    sharp corners, fontupper=\color{white}\bfseries},
  sharp corners=south, boxrule=1pt,
  top=5mm, bottom=4mm,
}{lem}

% COROLLARY
\newtcbtheorem[use counter=thmcounter]{corollary}{Corollary}{
  enhanced, breakable,
  colback=CorolBG, colframe=CorolBorder,
  arc=4pt,
  fonttitle=\bfseries\color{white},
  attach boxed title to top left={yshift=-\tcboxedtitleheight/2, xshift=8pt},
  boxed title style={colback=CorolBorder, colframe=CorolBorder,
    sharp corners, fontupper=\color{white}\bfseries},
  sharp corners=south, boxrule=1pt,
  top=5mm, bottom=4mm,
}{cor}

% DEFINITION
\newtcbtheorem[use counter=thmcounter]{definition}{Definition}{
  enhanced, breakable,
  colback=DefBG, colframe=DefBorder,
  arc=4pt,
  fonttitle=\bfseries\color{white},
  attach boxed title to top left={yshift=-\tcboxedtitleheight/2, xshift=8pt},
  boxed title style={colback=DefBorder, colframe=DefBorder,
    sharp corners, fontupper=\color{white}\bfseries},
  sharp corners=south, boxrule=1pt,
  top=5mm, bottom=4mm,
}{def}

% EXAMPLE
\newtcbtheorem[use counter=thmcounter]{example}{Example}{
  enhanced, breakable,
  colback=ExBG, colframe=ExBorder,
  arc=4pt,
  fonttitle=\bfseries\color{white},
  attach boxed title to top left={yshift=-\tcboxedtitleheight/2, xshift=8pt},
  boxed title style={colback=ExBorder, colframe=ExBorder,
    sharp corners, fontupper=\color{white}\bfseries},
  sharp corners=south, boxrule=0.8pt,
  top=5mm, bottom=4mm,
}{ex}

% REMARK  (lighter, no box title)
\newtcbtheorem[use counter=thmcounter]{remark}{Remark}{
  enhanced, breakable,
  colback=RemBG, colframe=RemBorder,
  arc=4pt,
  fonttitle=\bfseries\color{white},
  attach boxed title to top left={yshift=-\tcboxedtitleheight/2, xshift=8pt},
  boxed title style={colback=RemBorder, colframe=RemBorder,
    sharp corners, fontupper=\color{white}\bfseries},
  sharp corners=south, boxrule=0.6pt,
  top=5mm, bottom=4mm,
}{rem}

% PROBLEM
\newtcbtheorem[use counter=thmcounter]{problem}{Problem}{
  enhanced, breakable,
  colback=ProblemBG, colframe=ProblemBorder,
  arc=4pt,
  fonttitle=\bfseries\color{white},
  attach boxed title to top left={yshift=-\tcboxedtitleheight/2, xshift=8pt},
  boxed title style={colback=ProblemBorder, colframe=ProblemBorder,
    sharp corners, fontupper=\color{white}\bfseries},
  sharp corners=south, boxrule=1pt,
  top=5mm, bottom=4mm,
}{prob}

% ── PROOF style (white background, left gray bar, italic) ───
% Matches the screenshot: "Proof of Theorem X.X" header, plain body
\tcolorboxenvironment{proof}{
  blanker, breakable,
  left=4mm,
  before skip=6pt, after skip=10pt,
  borderline west={2pt}{0pt}{gray!60},
}
% Override proof header style
\renewcommand{\proofname}{\normalfont\color{gray}\small\scshape Proof}

% ── Code listings (listings + tcolorbox) ────────────────────
\usepackage{listings}

\lstdefinestyle{CppStyle}{
  language=C++,
  backgroundcolor=\color{CodeBG},
  basicstyle=\ttfamily\small\color{white},
  keywordstyle=\color{CodeKey}\bfseries,
  stringstyle=\color{CodeStr},
  commentstyle=\color{CodeCom}\itshape,
  numberstyle=\tiny\color{CodeNum},
  numbers=left, numbersep=8pt,
  stepnumber=1,
  breaklines=true, breakatwhitespace=false,
  tabsize=4,
  showstringspaces=false,
  morekeywords={auto, nullptr, constexpr, override, final,
                uint8_t, uint16_t, uint32_t, uint64_t,
                int8_t, int16_t, int32_t, int64_t},
  frame=none,
}

\lstdefinestyle{PythonStyle}{
  language=Python,
  backgroundcolor=\color{CodeBG},
  basicstyle=\ttfamily\small\color{white},
  keywordstyle=\color{CodeKey}\bfseries,
  stringstyle=\color{CodeStr},
  commentstyle=\color{CodeCom}\itshape,
  numberstyle=\tiny\color{CodeNum},
  numbers=left, numbersep=8pt,
  stepnumber=1,
  breaklines=true, breakatwhitespace=false,
  tabsize=4,
  showstringspaces=false,
  morekeywords={self, None, True, False, as, with, yield,
                lambda, import, from, pass, assert, raise,
                except, finally, elif},
  frame=none,
}

% Wrap lstlisting in a tcolorbox for rounded corners + title
\newtcblisting{cppcode}[1][]{
  listing only, listing style=CppStyle,
  colback=CodeBG, colframe=CodeBG!80!black,
  fonttitle=\ttfamily\small\color{white},
  title=\faCode\ C++,            % remove \faCode if not using fontawesome
  sharp corners, boxrule=0.5pt,
  #1
}

\newtcblisting{pycode}[1][]{
  listing only, listing style=PythonStyle,
  colback=CodeBG, colframe=CodeBG!80!black,
  fonttitle=\ttfamily\small\color{white},
  title=Python,
  sharp corners, boxrule=0.5pt,
  #1
}

% Fallback plain code block (language-agnostic)
\newtcblisting{code}[2][]{
  listing only,
  listing options={
    style=PythonStyle,
    language=#2,
  },
  colback=CodeBG, colframe=CodeBG!80!black,
  fonttitle=\ttfamily\small\color{white},
  title=#2,
  sharp corners, boxrule=0.5pt,
  #1
}

% ── Plots and Geometry (pgfplots + tikz) ────────────────────
\usepackage{tikz}
\usepackage{pgfplots}
\pgfplotsset{compat=1.18}
\usetikzlibrary{
  arrows.meta,         % Better arrow tips
  calc,                % Coordinate calculations
  shapes.geometric,    % circles, rectangles, etc.
  positioning,         % node placement
  decorations.markings,% marks on paths
  intersections,       % path intersections
  patterns,            % hatching
  angles, quotes,      % angle labeling
  3d,                  % 3D axes
}

% Convenience: TikZ picture wrapper with caption
\usepackage{float}
\usepackage{caption}
\usepackage{subcaption}

% ── Figures / Images ────────────────────────────────────────
\usepackage{graphicx}
\usepackage{wrapfig}   % Text-wrapped figures

% Figure box style (tcolorbox)
\newtcolorbox{figurebox}[2][]{
  enhanced, breakable,
  colback=white, colframe=gray!50,
  title=#2,
  fonttitle=\bfseries\small,
  boxed title style={colback=gray!20, colframe=gray!50},
  attach boxed title to top left={yshift=-\tcboxedtitleheight/2, xshift=8pt},
  sharp corners, boxrule=0.8pt,
  #1
}

% ── Tables ──────────────────────────────────────────────────
\usepackage{booktabs}   % \toprule, \midrule, \bottomrule
\usepackage{array}
\usepackage{multirow}
\usepackage{longtable}  % Tables spanning multiple pages

% ── Hyperlinks & PDF metadata ───────────────────────────────
\usepackage{hyperref}
\hypersetup{
  colorlinks = true,
  linkcolor  = LinkColor,
  citecolor  = LinkColor,
  urlcolor   = LinkColor,
  pdftitle   = {Lecture Notes},
  pdfauthor  = {Your Name},
}

% ── Table of Contents style ─────────────────────────────────
\usepackage{titletoc}
\usepackage{titlesec}

% Section numbering depth
\setcounter{tocdepth}{2}
\setcounter{secnumdepth}{3}

% Section title styling
\titleformat{\section}
  {\large\bfseries}
  {\thesection}{1em}{}[\titlerule]

\titleformat{\subsection}
  {\normalsize\bfseries}
  {\thesubsection}{1em}{}

% ── Bibliography ────────────────────────────────────────────
\usepackage[backend=biber, style=alphabetic]{biblatex}
% \addbibresource{references.bib}   % uncomment and add your .bib

% ── Miscellaneous Useful Packages ───────────────────────────
\usepackage{enumitem}          % Customise lists
\usepackage{todonotes}         % \todo{} margin notes
\usepackage{cleveref}          % \cref{} smart cross-refs
\usepackage{siunitx}           % SI units: \SI{9.8}{\metre\per\second}
\usepackage{physics}           % \dd, \grad, \curl, \laplacian …
\usepackage{mathrsfs}          % \mathscr{F} script font

% ── Convenience Math Macros ─────────────────────────────────
\newcommand{\R}{\mathbb{R}}
\newcommand{\N}{\mathbb{N}}
\newcommand{\Z}{\mathbb{Z}}
\newcommand{\Q}{\mathbb{Q}}
\newcommand{\C}{\mathbb{C}}
\newcommand{\E}{\mathbb{E}}
\newcommand{\Prob}{\mathbb{P}}
\newcommand{\one}{\mathbbm{1}}          % indicator
\newcommand{\eps}{\varepsilon}
% \newcommand{\norm}[1]{\left\lVert#1\right\rVert}
% \newcommand{\abs}[1]{\left\lvert#1\right\rvert}
\newcommand{\inner}[2]{\left\langle#1,\,#2\right\rangle}
\newcommand{\set}[1]{\left\{#1\right\}}
\newcommand{\floor}[1]{\left\lfloor#1\right\rfloor}
\newcommand{\ceil}[1]{\left\lceil#1\right\rceil}
\DeclareMathOperator*{\argmin}{arg\,min}
\DeclareMathOperator*{\argmax}{arg\,max}
\DeclareMathOperator{\sgn}{sgn}
% \DeclareMathOperator{\tr}{tr}
% \DeclareMathOperator{\rank}{rank}
\DeclareMathOperator{\diag}{diag}
\DeclareMathOperator{\supp}{supp}

% ── Header / Footer ─────────────────────────────────────────
\usepackage{fancyhdr}
\pagestyle{fancy}
\fancyhf{}
\fancyhead[L]{\small\textit{\CourseName}}
\fancyfoot[C]{\thepage}
\renewcommand{\headrulewidth}{0.4pt}

% ── Metadata Commands (fill these per lecture) ──────────────
\newcommand{\CourseName}{Linear Algebra II}
\newcommand{\CourseCode}{FMA\_2F001\_EP}
\newcommand{\University}{Ecole Polytechnique}
\newcommand{\Lecturer}{Prof. Guillaume Dubach}
\newcommand{\StudentName}{Nikolay Kharovskiy}
\newcommand{\Semester}{Fall 2026}

% ============================================================
%   DOCUMENT BEGIN
% ============================================================
\begin{document}

% ── Title Page ──────────────────────────────────────────────
\begin{titlepage}
  \centering
  \vspace*{2cm}

  {\Huge\bfseries \CourseName \par}
  \vspace{0.5cm}
  {\large \CourseCode \par}

  \vspace{1.5cm}
  \hrule height 1pt
  \vspace{0.4cm}
  {\Large\bfseries Lecture Notes \par}
  \vspace{0.4cm}
  \hrule height 1pt

  \vspace{2cm}

  \begin{tabular}{rl}
    \textbf{University:} & \University \\[4pt]
    \textbf{Lecturer:}   & \Lecturer   \\[4pt]
    \textbf{Student:}    & \StudentName \\[4pt]
    \textbf{Semester:}   & \Semester   \\[4pt]
  \end{tabular}

  \vfill

  {\small These are personal lecture notes and may contain errors.
         For the official course material consult the lecturer.}
\end{titlepage}

% ── Table of Contents ───────────────────────────────────────
\tableofcontents
\newpage

% ============================================================
%  LECTURE 1
% ============================================================
\section{Class 1: Basic Arithmetic}

The arithmetic fact $1+1=2$ was shown by Russel and Whitehead. The proof operates using Peano's axioms for the natural numbers $\N$.

\begin{theorem}{Peano's Axioms for $\N$}{}
\begin{itemize}
    \item \textbf{Axiom 1:} $0 \in \N$.
    \item \textbf{Axiom 2:} There exists a successor function $S: \N \rightarrow \N$.
\end{itemize}
From these, we can define $1=S(0)$ and $2=S(S(0))$.
\end{theorem}

\begin{definition}{Addition}{}
Addition is defined by the following rules:
\begin{enumerate}
    \item $\forall m \in \N: m+0 = m$ (Add1).
    \item $\forall m, n \in \N: m+S(n) = S(m+n)$ (Add2).
\end{enumerate}
\end{definition}

\begin{example}{Computing $1+1$}{}
$1+1 = 1+S(0) \stackrel{\text{(Add2)}}{=} S(1+0) \stackrel{\text{(Add1)}}{=} S(1) = 2$.
\end{example}

\begin{theorem}{Peano's Axiom 9 (Induction)}{}
Let $P$ be a unary predicate. To prove $\forall m \in \N, P(m)$ is true, it is enough to prove:
\begin{itemize}
    \item $P(0)$ is true.
    \item $\forall m: P(m) \implies P(S(m))$.
\end{itemize}
\end{theorem}

\begin{theorem}{Well Ordering Principle (WOP)}{}
Every non-empty subset of $\N$ has a minimum.
\end{theorem}

\begin{proposition}{WOP implies Induction}{}
The Well Ordering Principle can be used to prove the principle of Induction.
\end{proposition}
\begin{proof}
Let $S = \{m \mid P(m) \text{ is not true}\}$. If $P(m)$ is not true for every $m$, then $S$ is non-empty. We apply WOP to $S$, meaning there is a minimum element $m_0 = \min S$. But $m_0 \neq 0$, so $m_0 = S(m)$ for some $m$. This creates a contradiction.
\end{proof}

\subsection{Relations on a Set}

\begin{definition}{Relations and Equivalence}{}
Let $S$ be a set with elements $x, y \in S$. A relation on $S$ is defined as a subset of $S \times S$.
An equivalence relation must satisfy all three of the following properties:
\begin{itemize}
    \item \textbf{Reflexivity:} $\forall x \in S, x \sim x$.
    \item \textbf{Symmetry:} $\forall x, y \in S, x \sim y \iff y \sim x$.
    \item \textbf{Transitivity:} $\forall x, y, z \in S, x \sim y \text{ and } y \sim z \implies x \sim z$.
\end{itemize}
\end{definition}

\begin{example}{Equivalence Relations}{}
``Same Parity'' and ``Same remainder mod $n$'' are equivalence relations.
\end{example}

\begin{remark}{Partitions}{}
An equivalence relation is the same as a partition of $S$.
\end{remark}

\begin{definition}{Modular Classes}{}
For every $n \in \N$, $a \equiv b \pmod{n}$ defines an equivalence relation. The set of equivalence classes for this relation is denoted $\Z/n\Z$.
\end{definition}

\subsection{Modular Arithmetic and Divisibility}

\begin{theorem}{Chinese Remainder Theorem (Sunzi)}{}
If $m$ and $n$ are coprime, there is a bijection between $\Z/mn\Z$ and $\Z/m\Z \times \Z/n\Z$.
\end{theorem}

\begin{example}{Applying the Chinese Remainder Theorem}{}
For $n=2$ and $m=3$, we can map between $\Z/6\Z$ and $\Z/2\Z \times \Z/3\Z$.
If we want to find all $m \in \Z$ such that $m \equiv 2 \pmod{5}$ and $m \equiv 3 \pmod{7}$, the solution is the set $\{14 + 35k \mid k \in \Z\}$.
\end{example}

\begin{theorem}{Euclidean Division}{}
$\forall a, \forall b \neq 0, \exists (q, r)$ such that $a = bq + r$ and $0 \leq r < \vert{}b\vert{}$.
\end{theorem}

\begin{theorem}{Bézout's Theorem}{}
$\exists (u, v) \in \Z^2$ such that $au + bv = \gcd(a,b)$.
\end{theorem}

\begin{example}{Euclid's Algorithm}{}
Applying the algorithm to find coefficients for $5$ and $7$:
\begin{itemize}
    \item $7 = 1 \times 5 + 2$.
    \item $5 = 2 \times 2 + \gcd(5,7)$.
    \item $1 = 5 - 2 \times 2 = 5 - 2 \times (7-5) = 3 \times 5 - 2 \times 7$.
\end{itemize}
We have found that $3 \times 5 \rightarrow (0,1)$ and $-2 \times 7 \rightarrow (1,0)$. If we want to find $m$ such that $m \equiv 2 \pmod{5}$ and $m \equiv 3 \pmod{7}$, we compute $(3 \times 5) \times 3 + (-2 \times 7) \times 2$.
\end{example}

\subsection{Prime Numbers}

\begin{definition}{Prime Number}{}
A prime number is a number $p \in \N$ such that $p \geq 2$ and $p$ has only divisors $1$ and $p$.
\end{definition}

\begin{theorem}{Infinitude of Primes}{}
There are infinitely many primes.
\end{theorem}

\begin{theorem}{Tao's Theorem}{}
At least one of the following sets of primes $(p, q)$ is infinite:
\begin{itemize}
    \item \textbf{Twin primes:} $\vert{}p-q\vert{} = 2$.
    \item \textbf{Cousin primes:} $\vert{}p-q\vert{} = 4$.
    \item \textbf{Sexy primes:} $\vert{}p-q\vert{} = 6$.
\end{itemize}
\end{theorem}

\begin{remark}{Goldbach Conjecture}{}
The Goldbach conjecture posits that every even number is equal to $p+q$ with $p$ and $q$ primes.
\end{remark}

% ============================================================
%  LECTURE 2
% ============================================================
\section{Lecture 2: Groups}

\begin{definition}{Group}{}
A group is a set $G$ equipped with a binary operation (or group law) $*$, denoted $(G, *)$, that satisfies the following properties:
\begin{enumerate}
    \item \textbf{Associativity:} $\forall a, b, c \in G, (a * b) * c = a * (b * c)$.
    \item \textbf{Neutral Element:} $\exists e \in G$ such that $\forall a \in G, a * e = e * a = a$.
    \item \textbf{Inverse Element:} $\forall a \in G, \exists b \in G$ such that $a * b = b * a = e$.
\end{enumerate}
\textbf{Optional:} If a group also satisfies commutativity (i.e., $\forall a, b \in G, a * b = b * a$), it is called an \textbf{Abelian group}.
\end{definition}

\begin{example}{Groups in Modular Arithmetic}{}
\begin{itemize}
    \item The set $\Z/n\Z$ with addition is a group. The neutral element is $[0]_n$ and the inverse of $[k]_n$ is $[-k]_n$.
    \item If $p$ is prime, the non-zero elements modulo $p$ under multiplication form a group. To find the inverse of $k$ modulo $p$, one can use Bézout's identity on $k$ and $p$. Because $p$ is prime, there exist integers $u$ and $v$ such that $uk + vp = 1$. Modulo $p$, this simplifies to $uk \equiv 1 \pmod{p}$, which implies that $u$ is the multiplicative inverse $k^{-1}$.
\end{itemize}
\end{example}

\begin{definition}{Powers and Cyclic Groups}{}
For an element $g \in G$, its powers are defined as follows:
\begin{itemize}
    \item For $n > 0$, $g^n = g * g * \dots * g$ ($n$ times).
    \item For $n = 0$, $g^0 = e$.
    \item For $n < 0$, $g^n = g^{-1} * g^{-1} * \dots * g^{-1}$ ($|n|$ times).
\end{itemize}
A finite group $G$ is called \textbf{cyclic} if it is generated by a single element $g$, meaning $G = \langle g \rangle = \{ g^m \mid m \in \Z \}$.
\end{definition}

\begin{example}{Common Groups}{}
\begin{itemize}
    \item \textbf{Dihedral groups} $D_k$: the group of symmetries of a regular $k$-gon.
    \item \textbf{Symmetric groups} $\Sigma_n$ (or $S_n$): the group of permutations of $n$ elements.
    \item \textbf{Alternating group} $A_n$: the subgroup of $\Sigma_n$ containing all even permutations, i.e., $A_n = \{ \sigma \in \Sigma_n \mid \varepsilon(\sigma) = 1 \}$.
    \item The group of symmetries of polyhedra, such as the Rhombic Triacontahedron.
\end{itemize}
\end{example}

\begin{definition}{Subgroups and Morphisms}{}
\begin{itemize}
    \item A \textbf{subgroup} is a subset of $G$ that forms a group under the same group law.
    \item A \textbf{morphism} (or homomorphism) is a map $\varphi$ between groups that preserves the group structure, e.g., $\varphi(a * b) = \varphi(a) * \varphi(b)$.
    \item An \textbf{isomorphism} is a bijective group morphism.
\end{itemize}
\end{definition}

\begin{definition}{Order of an Element}{}
The order of an element $g \in G$ is the smallest positive integer $m$ such that $g^m = e$. It is denoted $\omega(g) = \min \{ m \ge 1 \mid g^m = e \}$.
\end{definition}

\begin{theorem}{Weak Lagrange Theorem}{}
Let $G$ be a finite Abelian group, and let $g \in G$. Then the order of $g$ divides the size of the group, i.e., $\omega(g) \mid \#G$.
\end{theorem}
\begin{proof}
Let $\pi = \prod_{h \in G} h$ be the product of all elements in $G$. Consider the map $\varphi_g : G \to G$ defined by $h \mapsto g * h$. This map is a bijection, as it has an inverse $\varphi_{g^{-1}}$. Because $G$ is Abelian, we can rearrange the product:
$$ \pi = \prod_{h \in G} h = \prod_{h \in G} \varphi_g(h) = \prod_{h \in G} (g * h) $$
Factoring out $g$ for each of the $\#G$ elements, we get $\pi = g^{\#G} \cdot \prod_{h \in G} h = g^{\#G} \cdot \pi$. By canceling $\pi$, we find $g^{\#G} = e$.

We can write $\#G = q \cdot \omega(g) + r$ using Euclidean division, with $0 \le r < \omega(g)$. Substituting this into our equation gives:
$$ e = g^{\#G} = g^{q \omega(g) + r} = (g^{\omega(g)})^q \cdot g^r = e^q \cdot g^r = g^r $$
Since $g^r = e$ and $0 \le r < \omega(g)$, the minimality of the order implies $r = 0$. Therefore, $\omega(g)$ divides $\#G$.
\end{proof}

\begin{theorem}{Strong Lagrange Theorem}{}
If $G$ is a finite group and $H \subset G$ is a subgroup, then $\#H$ divides $\#G$.
\end{theorem}
\begin{proof}
We say two elements $x, y \in G$ are equivalent modulo $H$ ($x \sim y$) if $y^{-1}x \in H$. This condition is equivalent to $x \in yH$, where $yH$ represents a left coset. The left cosets form a partition of the group $G$, so $G = \bigsqcup_{y_i \in I} y_iH$ for some index set $I$ of representatives.

Every coset $y_iH$ has exactly the same size, which is equal to $\#H$. Thus, we can conclude that the total number of elements in $G$ is given by $\#G = \#H \times \#I$. This immediately shows that $\#H \mid \#G$.
\end{proof}

\begin{theorem}{Cauchy's Theorem}{}
Let $G$ be a finite group and let $p$ be a prime number such that $p \mid \#G$. Then there exists an element $g \in G$ such that $\omega(g) = p$.
\end{theorem}
\begin{proof}
Consider the set $S_p = \{ (g_1, g_2, \dots, g_p) \mid g_i \in G, g_1 g_2 \dots g_p = e \}$. 
We define an equivalence relation on $S_p$ using cyclic shifts: two tuples are equivalent if one is a cyclic permutation of the other (i.e., shifting the sequence by some offset $k$).
Notice that the size of $S_p$ is $(\#G)^{p-1}$, because the first $p-1$ elements can be chosen freely, and the final element $g_p$ is uniquely determined as $(g_1 \dots g_{p-1})^{-1}$.

Under this equivalence relation, every equivalence class contains either $1$ or $p$ elements. The tuple $(e, \dots, e)$ forms an equivalence class of size $1$.
We can partition $S_p$ into these classes, yielding the equation $\#S_p = N_1 + p N_p$, where $N_1$ is the number of classes of size $1$, and $N_p$ is the number of classes of size $p$.

Since $p \mid \#G$, we know $p \mid (\#G)^{p-1}$, so $p \mid \#S_p$. Thus, $p \mid (N_1 + p N_p)$, which implies $p \mid N_1$.
We already know $N_1 \ge 1$ because of the class containing $(e, \dots, e)$. Because $p \mid N_1$, we must have $N_1 \ge p \ge 2$.

This means there is at least one other class of size $1$, which corresponds to a tuple $(g, g, \dots, g) \in S_p$ for some $g \neq e$. Since this tuple is in $S_p$, it must satisfy $g^p = e$. Because $p$ is prime and $g \neq e$, the order of $g$ is exactly $p$, so $\omega(g) = p$.
\end{proof}

% ============================================================
%  LECTURE 3
% ============================================================
\section{Lecture 3: Rings}

\begin{definition}{Unital Ring}{}
A unital ring is a set $R$ equipped with two binary operations, denoted $(R, +, \times)$, satisfying the following properties:
\begin{enumerate}
    \item $(R, +)$ is an Abelian group, with a neutral element denoted $0$ or $0_R$. %[cite: 16]
    \item Multiplication "$\times$" is associative: $\forall x, y, z \in R, x(yz) = (xy)z$. %[cite: 16]
    \item Multiplication "$\times$" has a neutral element, denoted $1$ or $1_R$. %[cite: 16]
    \item Distributivity of multiplication over addition: $\forall a, b, c \in R, a(b+c) = ab + ac$ and $(b+c)a = ba + ca$. %[cite: 16]
\end{enumerate}
\textbf{Note:} $(R, \times)$ and $(R \setminus \{0\}, \times)$ are not necessarily groups, because elements do not generally have multiplicative inverses. %[cite: 16]

If the multiplication operation is also commutative, then $R$ is called a \textbf{commutative ring}. %[cite: 16]
\end{definition}

\begin{definition}{Ring Morphism}{}
A map $\varphi : (R_1, +, \times) \to (R_2, +, \times)$ is a ring morphism if:
\begin{enumerate}
    \item $\varphi : (R_1, +) \to (R_2, +)$ is a group morphism. %[cite: 16, 17]
    \item $\varphi(1_{R_1}) = 1_{R_2}$. %[cite: 17]
    \item $\forall a, b \in R_1, \varphi(ab) = \varphi(a)\varphi(b)$. %[cite: 17]
\end{enumerate}
\end{definition}

\begin{theorem}{The Ring of Integers as an Initial Object}{}
The ring $(\Z, +, \times)$ is the paradigmatic example of a ring, and it serves as the initial object in the category of unital rings. %[cite: 17]
This means that for any ring $R$, there exists a unique ring morphism $\psi : \Z \to R$. %[cite: 17]
Therefore, $\Z$ is the unique initial object up to ring isomorphism. %[cite: 17]
\end{theorem}
\begin{proof}
By "abstract nonsense", let $R_0$ be another initial object. %[cite: 18]
Then there exist unique morphisms $\psi_1 : \Z \to R_0$ and $\psi_2 : R_0 \to \Z$. %[cite: 18]
The composition $\psi_2 \circ \psi_1 : \Z \to \Z$ must be the identity morphism (since $\Z$ is initial), which implies $\psi_1 = \psi_2^{-1}$, establishing a ring isomorphism between $\Z$ and $R_0$. %[cite: 18]
\end{proof}

\subsection{Subrings and Ideals}

\begin{definition}{Subring}{}
A subring is a subset that forms a ring under the same (induced) algebraic structure. %[cite: 18]
For example, we have the chain of subrings: $\Z \subset \Q \subset \R \subset \C \subset \mathbb{H}$. %[cite: 18]
\end{definition}

\begin{remark}{Ideals vs. Subrings}{}
For groups, the kernel of a group morphism is a subgroup. %[cite: 18]
However, for rings, the kernel is generally \textit{not} a subring (it usually does not contain $1_R$). %[cite: 18, 19]
The correct notion to capture the behavior of kernels in rings is an \textbf{ideal}. %[cite: 19]
\end{remark}

\begin{definition}{Ideal}{}
An ideal $I \subset R$ is a subgroup of $(R, +)$ that is stable under multiplication by any element of $R$. %[cite: 19]
Equivalently: $\forall a, b \in R$ and $\forall x, y \in I$, we have $ax + by \in I$ (analogous to a linear combination in a vector space). %[cite: 19]
\end{definition}

\begin{proposition}{Kernels are Ideals}{}
Let $\varphi : R_1 \to R_2$ be a ring morphism. %[cite: 19]
Then $\ker(\varphi) = \{ x \in R_1 \mid \varphi(x) = 0 \}$ is an ideal of $R_1$ (but not a subring!). %[cite: 19]
\end{proposition}
\begin{proof}
Let $a, b \in R_1$ and $x, y \in \ker(\varphi)$. %[cite: 19]
We want to prove that $ax + by \in \ker(\varphi)$, which means $\varphi(ax + by) = 0$. %[cite: 20]
Using the properties of ring morphisms:
$$ \varphi(ax + by) = \varphi(ax) + \varphi(by) = \varphi(a)\varphi(x) + \varphi(b)\varphi(y) $$ %[cite: 20]
Since $x, y \in \ker(\varphi)$, we know $\varphi(x) = 0$ and $\varphi(y) = 0$. %[cite: 20]
Thus, $\varphi(a)\cdot 0 + \varphi(b)\cdot 0 = 0$. %[cite: 20]
\end{proof}

\begin{definition}{Principal Ideal}{}
In a commutative ring $R$, ideals of the form $(x) = xR = \{xy \mid y \in R\}$ are called \textbf{principal ideals}. %[cite: 20]
For example, all ideals of $\Z$ are principal (of the form $m\Z$). %[cite: 20]
\end{definition}

\subsection{Quotient Rings}

\begin{theorem}{Quotient Ring Construction}{}
For any commutative ring $R$ and any ideal $I \subset R$, we can construct a quotient space $R/I$ by defining an equivalence relation $r_1 \sim r_2 \iff r_1 - r_2 \in I$. %[cite: 20, 21]
The equivalence classes (or cosets) are denoted $[r] = r + I$. %[cite: 21]

The set of these classes, $R/I$, inherits a natural ring structure given by:
\begin{itemize}
    \item $[r_1] + [r_2] = [r_1 + r_2]$ %[cite: 21]
    \item $[r_1] \times [r_2] = [r_1 \times r_2]$ %[cite: 21]
\end{itemize}
\end{theorem}

\begin{example}{Modular Arithmetic and CRT}{}
Let $R = \Z$ and $I = m\Z$. %[cite: 22]
This construction gives the usual quotient ring $\Z/m\Z$. %[cite: 22]
If $m$ and $n$ are coprime, the Chinese Remainder Theorem establishes a ring isomorphism: $\Z/mn\Z \cong \Z/m\Z \times \Z/n\Z$. %[cite: 22]
\end{example}

\subsection{Taxonomy of Rings}

\begin{definition}{Types of Rings}{}
\begin{itemize}
    \item \textbf{Integral Domain:} A commutative ring where $xy = 0 \implies x = 0 \text{ or } y = 0$ (i.e., it has no zero divisors). %[cite: 22]
    \item \textbf{Principal Ideal Domain (PID):} An integral domain where every ideal is a principal ideal. %[cite: 22]
    \item \textbf{Division Ring:} A ring where $(R \setminus \{0\}, \times)$ is a group (i.e., every non-zero element has a multiplicative inverse). %[cite: 23]
    \item \textbf{Field:} A commutative division ring (i.e., $(R \setminus \{0\}, \times)$ is an Abelian group). %[cite: 23]
\end{itemize}
\end{definition}

\begin{remark}{Examples and Hierarchy}{}
We have the following strict inclusions for commutative rings:
$$ \text{Commutative Rings} \supset \text{Integral Domains} \supset \text{PIDs} \supset \text{Fields} $$ %[cite: 23]
\begin{itemize}
    \item \textbf{Non-commutative Rings:} $M_2(\R), M_3(\R)$. %[cite: 23]
    \item \textbf{Division Rings (not fields):} The quaternions $\mathbb{H}$. %[cite: 23]
    \item \textbf{Fields:} $\R, \Q, \Z/p\Z$ (where $p$ is prime). %[cite: 23]
\end{itemize}
\end{remark}

\subsection{Special Ideals}

\begin{definition}{Prime and Maximal Ideals}{}
Let $I \subsetneq R$ be a proper ideal in a commutative ring $R$.
\begin{itemize}
    \item $I$ is a \textbf{prime ideal} if $xy \in I \implies x \in I \text{ or } y \in I$. %[cite: 23]
    For example, $2\Z$ is a prime ideal, but $4\Z$ is not. %[cite: 23]
    \item $I$ is a \textbf{maximal ideal} if for any ideal $J$ such that $I \subseteq J \subsetneq R$, it must be that $I = J$. %[cite: 24]
\end{itemize}
\end{definition}

\begin{theorem}{Quotients by Prime and Maximal Ideals}{}
Let $R$ be a commutative ring and $I$ be a proper ideal. %[cite: 24]
\begin{itemize}
    \item $I$ is a maximal ideal $\iff (R/I, +, \times)$ is a field. %[cite: 24]
    \item $I$ is a prime ideal $\iff (R/I, +, \times)$ is an integral domain. %[cite: 24]
\end{itemize}
Because every field is an integral domain, this immediately implies that every maximal ideal is also a prime ideal. %[cite: 24]
\end{theorem}

% ============================================================
%  LECTURE 4
% ============================================================
\section{Lecture 4: Fields and Polynomials}

\begin{definition}{Field}{}
A set $K$ equipped with two operations $(K, +, \times)$ is a field if:
\begin{enumerate}
    \item It is a commutative ring. 
    \item It is a division ring (i.e., every non-zero element has a multiplicative inverse, so $\forall x \neq 0, \exists y$ such that $xy = 1$). 
\end{enumerate}
\end{definition}

\begin{remark}{Properties of Fields}{}
\begin{itemize}
    \item The non-zero elements form a group under multiplication, denoted $(K^*, \times)$. 
    \item Every field $(K, +, \times)$ is an integral domain (since $xy = 0 \iff x = 0 \text{ or } y = 0$). 
\end{itemize}
\end{remark}

\begin{definition}{Field Extensions}{}
A \textbf{subfield} is a subset of a field that is itself a field under the induced operations. 
If $K \subset L$ is a subfield, we usually refer to $L$ as a \textbf{field extension} of $K$. 
Field extensions often arise to supply "missing" roots to equations.  
For example, $\C$ extends $\R$ to provide roots for $X^2 + 1 = 0$, and $\R$ extends $\Q$ to provide roots for polynomials like $X^2 - 2$, $X^2 - 3$, or $X^3 - 4$. 
\end{definition}

\subsection{Polynomial Rings}

\begin{definition}{The Polynomial Ring $K[X]$}{}
Let $K$ be a field. The set of polynomials with coefficients in $K$ is denoted $K[X]$. 
An element $P \in K[X]$ is of the form $P(X) = a_0 + a_1X + \dots + a_nX^n$, where $(a_0, \dots, a_n)$ is a finite sequence of elements in $K$. 
\begin{itemize}
    \item $(K[X], +, \times)$ is a commutative ring, but \textbf{not} a division ring. 
    \item If $a_n \neq 0$, then the \textbf{degree} of $P$ is $\deg(P) = n$. 
    \item By convention, the zero polynomial $P = 0$ has degree $\deg(0) = -\infty$. 
\end{itemize}
\end{definition}

\begin{proposition}{Degree and Integral Domain}{}
For any polynomials $P, Q \in K[X]$:
\begin{enumerate}
    \item $\deg(PQ) = \deg(P) + \deg(Q)$. 
    \item $K[X]$ is an integral domain. 
\end{enumerate}
\end{proposition}
\begin{proof}
If $PQ = 0$, then $\deg(PQ) = -\infty$. 
Therefore, $\deg(P) + \deg(Q) = -\infty$, which requires either $\deg(P) = -\infty$ or $\deg(Q) = -\infty$. 
Thus, $P = 0$ or $Q = 0$. 
\end{proof}

\begin{remark}{Polynomials as Functions}{}
Every polynomial $P \in K[X]$ can be identified with a polynomial function $K \to K$, mapping $\alpha \mapsto P(\alpha) = a_0 + a_1\alpha + \dots + a_n\alpha^n$. 
An element $\alpha \in K$ is called a \textbf{root} of $P$ if $P(\alpha) = 0$. 
\end{remark}

\subsection{Euclidean Division and Ideals in $K[X]$}

\begin{theorem}{Euclidean Division in $K[X]$}{}
For any polynomials $A, B \in K[X]$ with $B \neq 0$, there exist unique polynomials $Q, R \in K[X]$ such that:
$$ A(X) = B(X)Q(X) + R(X) \quad \text{with} \quad \deg(R) < \deg(B) $$ 
\end{theorem}

\begin{example}{Computing Euclidean Division}{}
Let $A(X) = X^2 + 3X + 1$ and $B(X) = 2X + 2$ in $\Q[X]$. 
We can rewrite $A$ as:
$$ X^2 + 3X + 1 = (2X + 2)\left(\frac{1}{2}X + 1\right) - 1 $$ 
Here, the quotient is $Q(X) = \frac{1}{2}X + 1$ and the remainder is $R(X) = -1$. 
\end{example}

\begin{theorem}{$K[X]$ is a Principal Ideal Domain (PID)}{}
Every ideal in $K[X]$ is principal. 
\end{theorem}
\begin{proof}
Let $I \subset K[X]$ be a non-zero ideal. We want to show that $I = (B) = B \cdot K[X]$ for some $B \in I$. 
If $I$ contains a non-zero constant polynomial, then $1 \in I \implies I = K[X]$, which is generated by the constant polynomial $1$. 

Otherwise, by the Well Ordering Principle, let $B \in I$ be a non-zero polynomial of smallest degree in $I$.
Let $A \in I$ be any polynomial. We perform Euclidean division of $A$ by $B$: 
$$ A = BQ + R \quad \text{where} \quad \deg(R) < \deg(B) $$ 
Since $A \in I$ and $B \in I$, and $I$ is an ideal, the remainder $R = A - BQ$ must also be in $I$. 
However, $\deg(R) < \deg(B)$, and $B$ was chosen to have the minimal degree among non-zero elements in $I$. 
This implies $R$ must be the zero polynomial ($R = 0$). 
Therefore, $A = BQ$, meaning every element in $I$ is a multiple of $B$. Thus, $I = (B)$. 
\end{proof}

\subsection{Roots and the Fundamental Theorem of Algebra}

\begin{proposition}{Roots correspond to Linear Factors}{}
If $\alpha \in K$ is a root of $P \in K[X]$, then $P(X) = (X - \alpha)Q(X)$ for some $Q \in K[X]$. 
\end{proposition}
\begin{proof}
Perform Euclidean division of $P(X)$ by $(X - \alpha)$:
$$ P(X) = (X - \alpha)Q(X) + R(X) $$ 
The remainder must satisfy $\deg(R) < \deg(X - \alpha) = 1$, which means $\deg(R) \le 0$. 
Thus, $R(X)$ is a constant polynomial $r$. 
Evaluating the polynomial at $\alpha$ gives $P(\alpha) = (\alpha - \alpha)Q(\alpha) + r = r$. 
Since $\alpha$ is a root, $P(\alpha) = 0$, so $r = 0$. 
Therefore, $P(X) = (X - \alpha)Q(X)$. 
\end{proof}

\begin{corollary}{Maximum Number of Roots}{}
A non-zero polynomial $P \in K[X]$ of degree $n$ has at most $n$ roots in $K$. 
\end{corollary}
\begin{proof}
For every root found, the polynomial can be factored by a degree $1$ factor, reducing the degree of the quotient by $1$. This process can be repeated at most $n$ times. 
\end{proof}

\begin{theorem}{Fundamental Theorem of Algebra}{}
\textit{(Referred to in lecture as an "insignificant analytic lemma")} \\
Every polynomial $P \in \C[X]$ with $\deg(P) \ge 1$ has at least one root in $\C$. 
\end{theorem}

\begin{corollary}{Factorization in $\C[X]$}{}
Every polynomial $P \in \C[X]$ of degree $n \ge 1$ has exactly $n$ roots in $\C$ (counted with multiplicity). 
It factors completely into linear terms:
$$ P(X) = c(X - z_1)(X - z_2)\dots(X - z_n) $$ 
where $c$ is the leading coefficient.
\end{corollary}

\begin{proof}[Proof Sketch of the Fundamental Theorem of Algebra]
Let $P(z) = a_0 + a_1z + \dots + a_nz^n$ with $a_n \neq 0$ and $n \ge 1$. 
If $a_0 = 0$, then $P(0) = 0$, so $z=0$ is a root. Assume $a_0 \neq 0$. 

Consider the image of a circle of radius $r$ under the polynomial map, parameterized by $z = re^{i\theta}$. 
For a very large radius $r$, the highest degree term dominates, meaning $P(re^{i\theta}) \approx a_n r^n e^{in\theta}$. 
This creates a large curve that winds around the origin $n$ times. 
As $r \to 0$, the curve shrinks continuously to the point $P(0) = a_0$, which is non-zero (so it no longer winds around the origin). 
By continuity, as the curve deforms from winding around the origin to a point away from the origin, the curve must cross $0$ at some radius $r$ and angle $\theta$. 
At this intersection, $P(re^{i\theta}) = 0$, giving a root. 
\end{proof}

\end{document}